\section{Complex cotangent powers and their resonant amplitudes}
\label{bhr:cxp:sec:complex}

Question R4 of the incoming \emph{Harmonic Parity and Resolvent
Identities} report \cite{bhr:Parity} asks for a joint generator for nonintegral powers of
the cotangent resolvent, with the moving endpoint amplitudes retained
before expansion.  We prove such a generator and its finite-part
conversion.  The primary result below covers every complex exponent
with positive real part and every positive integral crossing.  A new
Dirichlet-series coordinate is used explicitly; no finite reduction to
ordinary polylogarithms is asserted for an arbitrary complex exponent.
Nevertheless, its exponent jets have a finite multiple-polylogarithmic
description, and one special resolvent point has a Gamma evaluation at
every positive integral spectral order.

\subsection{Branches and a binomial Dirichlet transform}

Put
\[
 K(x)=\pi\cot(\pi x),\qquad R_z(x)=(K(x)-z)^{-1},
 \qquad z\in\mathbb C\setminus\mathbb R.
\]
For the chosen half-plane define
\[
 \epsilon=\operatorname{sgn}\Im z,\quad
 a=z+\epsilon\pi i,\quad A=-a^{-1},\quad
 q=\frac{z-\epsilon\pi i}{z+\epsilon\pi i},\quad
 \sigma=-2\epsilon\pi i.
\]
Then \(|q|<1\).  We take the continuous logarithm of \(K(x)-z\)
whose imaginary part lies in \((-\pi,0)\) if \(\epsilon=1\) and
in \((0,\pi)\) if \(\epsilon=-1\).  Consequently
\[
 R_z(x)^\lambda=\exp\{-\lambda\Log(K(x)-z)\}.
\]
Also \(\Log\sigma=\log(2\pi)-\epsilon\pi i/2\), and
\(\Log A\) is the principal logarithm.

With \(w=e^{2\epsilon\pi i x}\), elementary algebra gives
\begin{equation}\label{bhr:cxp:eq:disk}
 R_z(x)^\lambda
 =A^\lambda(1-w)^\lambda(1-qw)^{-\lambda}.
\end{equation}
Here the two logarithms on the right are their holomorphic
determinations in \(|w|<1\), continued to \(w\ne1\) on its boundary.
To check the constant in this identity, both sides divided by
\(x^\lambda\) tend to one as \(x\downarrow0\).  This fixes the phase
at the other endpoint as well.

Define
\begin{align}
 h_n(\lambda,q)
 &=[w^n](1-w)^\lambda(1-qw)^{-\lambda}\notag\\
 &=\sum_{k=0}^n
   \frac{(-\lambda)_k(\lambda)_{n-k}}{k!(n-k)!}\,q^{n-k},
 \label{bhr:cxp:eq:hn}\\
 \mathcal L_\lambda(u,q)
 &=\sum_{n=1}^{\infty}\frac{h_n(\lambda,q)}{n^u}.
 \label{bhr:cxp:eq:L}
\end{align}
The Pochhammer symbols have their polynomial meaning, including at
integral parameters.  The initial Dirichlet series converges normally
where \(\Re(u+\lambda)>0\), locally in \(|q|<1\).  One way to see
the needed coefficient estimate is to convolve the binomial
coefficients of \((1-w)^\lambda\) with the exponentially decaying
coefficients of \((1-qw)^{-\lambda}\); on compact parameter sets this
gives \(h_n=O(n^{-\Re\lambda-1})\), with a harmless uniform weakening
of the exponent at the boundary of a compact set.

\begin{theorem}[Joint continuation of the binomial transform]
\label{bhr:cxp:thm:L}
The function \(\mathcal L_\lambda(u,q)\) continues meromorphically
in \((u,\lambda)\in\mathbb C^2\), holomorphically in \(|q|<1\).
Its possible polar hypersurfaces are
\[
 u+\lambda=-j,\qquad j=0,1,2,\ldots.
\]
Away from them,
\[
 \mathcal L_\lambda(0,q)=-1,\qquad
 \mathcal L_\lambda(-N,q)=0\quad(N\ge1).
\]
These equalities concern the meromorphic restriction at the indicated
spectral order.  At an intersection with a polar hypersurface,
substitution and continuation need not commute.
\end{theorem}

\begin{proof}
For \(\Re u>0\) and \(\Re(u+\lambda)>0\), termwise integration
gives
\begin{equation}\label{bhr:cxp:eq:mellin}
 \Gamma(u)\mathcal L_\lambda(u,q)
 =\int_0^\infty t^{u-1}
 \left[
 \left(\frac{1-e^{-t}}{1-qe^{-t}}\right)^\lambda-1
 \right]\,dt.
\end{equation}
The positive real logarithm of \(1-e^{-t}\), and the holomorphic
logarithm of \(1-qe^{-t}\), specify its powers.  Set
\[
 \left(\frac{1-e^{-t}}{1-qe^{-t}}\right)^\lambda
   =t^\lambda\sum_{j\ge0}b_j(\lambda,q)t^j.
\]
The coefficients are entire in \(\lambda\), holomorphic in \(q\).
For any fixed \(N\), split the integral at one and subtract the
first \(N\) terms of this local expansion.  The result is
\begin{align}
 \Gamma(u)\mathcal L_\lambda(u,q)
 ={}&-\frac1u+
   \sum_{j=0}^{N-1}\frac{b_j(\lambda,q)}{u+\lambda+j}
   +E_N(u,\lambda,q),\label{bhr:cxp:eq:subtraction}
\end{align}
where \(E_N\) is holomorphic when \(\Re(u+\lambda)>-N\).
The tail at infinity is exponentially convergent on every compact
parameter set.  Increasing \(N\) proves the continuation.  Division
by \(\Gamma(u)\) removes the isolated pole \(-1/u\), with value
\(-1\) at \(u=0\), and gives zero at the other nonpositive
integers whenever no moving denominator vanishes.
\end{proof}

\subsection{The completed spectral generator}

Use the conventional normalization
\[
 \zeta(1-s,x)=-\frac1s+
    \sum_{m\ge0}\gamma_m(x)\frac{s^m}{m!},
 \qquad \gamma_0(x)=-\psi(x).
\]
For an integer \(p\ge0\), let
\[
 A_p(s)=\prod_{k=1}^p(s-k),\qquad A_0(s)=1.
\]
Argument differentiation satisfies
\[
 \partial_x^p\zeta(1-s,x)
   =A_p(s)\zeta(1+p-s,x).
\]
The notation \(\gamma_m^{(p)}\) below always means argument
differentiation.

\begin{theorem}[Complex-power cotangent generator]
\label{bhr:cxp:thm:generator}
In the ordinary-integral region
\(\Re\lambda>-1,\ \Re(s+\lambda)>p\), away from the Hurwitz
pole when \(p=0\),
\begin{equation}\label{bhr:cxp:eq:J}
 \mathcal J_p(s,\lambda,z):=
 \int_0^1\partial_x^p\zeta(1-s,x)R_z(x)^\lambda\,dx
 =A^\lambda\Gamma(s)\sigma^{p-s}
       \mathcal L_\lambda(s-p,q).
\end{equation}
This identity defines a joint meromorphic continuation to all
\((s,\lambda)\).  Besides the pole
\(-A^\lambda/s\) when \(p=0\), its only possible polar
hypersurfaces are
\[
 s+\lambda=p-j,\qquad j=0,1,\ldots.
\]
Define the completed family
\begin{equation}\label{bhr:cxp:eq:F}
 \mathcal F_p(s,\lambda,z)=
 \mathcal J_p(s,\lambda,z)
       +\mathbf1_{p=0}\frac{A^\lambda}{s}.
\end{equation}
For \(\Re\lambda>p\) the ordinary absolutely convergent moments
\begin{equation}\label{bhr:cxp:eq:M}
 \mathcal M_{m,p}(\lambda,z)=
 \int_0^1\gamma_m^{(p)}(x)R_z(x)^\lambda\,dx
 =m![s^m]\mathcal F_p(s,\lambda,z)
\end{equation}
therefore have meromorphic continuations in \(\lambda\), with
possible poles only at the integers \(\lambda\le p\), of order
at most \(m+1\).
\end{theorem}

\begin{proof}
First suppose \(\Re\lambda>0\) and \(\Re s>p+1\).
The Fourier series of the kernel in \eqref{bhr:cxp:eq:disk}
is absolutely convergent.  The classical Hurwitz Fourier identity
\cite[25.11.9]{bhr:DLMFZ} gives
\[
 \int_0^1\zeta(1-s,x)\,dx=0,\qquad
 \int_0^1\partial_x^p\zeta(1-s,x)
           e^{2\epsilon\pi i n x}\,dx
 =\Gamma(s)\sigma^{p-s}n^{p-s}.
\]
The constant kernel mode vanishes, and summing the other modes
proves \eqref{bhr:cxp:eq:J}.  The endpoint decomposition
\[
 \partial_x^p\zeta(1-s,x)
 =A_p(s)x^{s-p-1}
       +\partial_x^p\zeta(1-s,1+x)
\]
shows normal integrability on the larger region stated in the
theorem, so the identity extends there.

In \eqref{bhr:cxp:eq:subtraction}, put \(u=s-p\) and use
\(\Gamma(s)/\Gamma(s-p)=A_p(s)\).  The moving denominators
give exactly the indicated hypersurfaces.  The isolated term
\(-1/(s-p)\) is cancelled by \(A_p(s)\) for \(p\ge1\); for
\(p=0\) it gives residue \(-A^\lambda\).  This proves the
claimed completion.  Finally, the mean of the kernel is \(A^\lambda\)
for \(\Re\lambda>0\), by its constant Fourier coefficient.
Termwise Stieltjes expansion is valid when \(\Re\lambda>p\),
including all powers of the endpoint logarithm.  This proves
\eqref{bhr:cxp:eq:M} and its continuation.
\end{proof}

\begin{corollary}[Lowest Stieltjes index and shift identities]
\label{bhr:cxp:cor:shifts}
Away from the moving poles,
\begin{align*}
 \mathcal M_{0,0}(\lambda,z)
 &=A^\lambda\left(
   \gamma+\Log\sigma+
         \left.\partial_u\mathcal L_\lambda(u,q)\right|_{u=0}\right),\\
 \mathcal M_{0,p}(\lambda,z)
 &=A^\lambda\sigma^p
         \left.\partial_u\mathcal L_\lambda(u,q)\right|_{u=-p}
                                      \qquad(p\ge1).
\end{align*}
The continued moments satisfy the meromorphic identities
\begin{align}
 \partial_z\mathcal M_{m,p}(\lambda,z)
   &=\lambda\mathcal M_{m,p}(\lambda+1,z),\label{bhr:cxp:eq:zshift}\\
 \mathcal M_{m,p+1}(\lambda,z)
   &=-\lambda\bigl(
       \mathcal M_{m,p}(\lambda-1,z)
       +2z\mathcal M_{m,p}(\lambda,z)\notag\\
   &\hspace{39mm}
       +(z^2+\pi^2)\mathcal M_{m,p}(\lambda+1,z)\bigr).
       \label{bhr:cxp:eq:pshift}
\end{align}
\end{corollary}
\begin{proof}
The first two formulas follow by expanding \(\Gamma(s)\sigma^{-s}\)
and using Theorem~\ref{bhr:cxp:thm:L} at \(0,-p\).
The first shift follows from
\(\partial_zR_z^\lambda=\lambda R_z^{\lambda+1}\).
For the second, integrate by parts where \(\Re\lambda>p+1\),
so the boundary terms vanish, and use
\[
 R_z'=1+2zR_z+(z^2+\pi^2)R_z^2.
\]
Continuation proves the meromorphic statements.  At a resonant
integer, a pointwise finite-part version must use the conversion
in the next theorem before substituting.
\end{proof}

\subsection{All positive resonances and the two finite parts}

Finite parts in the next theorem use the original \(x\)-coordinate:
they retain the constant coefficient in the cutoff integral
\(\int_\eta^{1-\eta}\), after removing every nonintegrable power
and logarithm.  For a complex exponent of positive real part this
definition is unambiguous from the local power-log expansion.
It agrees with the ordinary integral whenever that integral converges.

The complete local amplitude is elementary.  Write
\begin{equation}\label{bhr:cxp:eq:d}
 D_j(\lambda,z)
 =[x^j]\bigl(\pi x\cot(\pi x)-zx\bigr)^{-\lambda}.
\end{equation}
Thus \(R_z(x)^\lambda
=x^\lambda\sum_{j\ge0}D_j(\lambda,z)x^j\) near zero, with
\[
 D_0=1,\qquad D_1=\lambda z,\qquad
 D_2=\frac{\lambda\pi^2}{3}
       +\frac{\lambda(\lambda+1)z^2}{2}.
\]
Every \(D_j\) is a polynomial in \(\lambda,z,\pi^2\) with rational
coefficients, and has degree at most \(j\) in \(\lambda\).

\begin{theorem}[Resonant amplitudes and finite-part conversion]
\label{bhr:cxp:thm:contacts}
For a positive integer \(r\le p\), put \(j=p-r\).
In a neighborhood of \((s,\lambda)=(0,r)\),
\begin{equation}\label{bhr:cxp:eq:amplitude}
 \mathcal F_p(s,\lambda,z)
  =\frac{A_p(s)D_j(\lambda,z)}{s+\lambda-r}
      +H_{p,r}(s,\lambda,z),
\end{equation}
where \(H_{p,r}\) is jointly holomorphic.  Hence, for every
positive integer \(r\),
\begin{align}
 &\operatorname{FP}_x\int_0^1
       \gamma_m^{(p)}(x)R_z(x)^r\,dx\notag\\
 &\quad =
 m![s^m]\left\{
   \mathcal F_p(s,r,z)
   -\mathbf1_{r\le p}\frac{A_p(s)D_{p-r}(r,z)}s
           \right\}. \label{bhr:cxp:eq:fixed-r}
\end{align}

There is also an explicit comparison with continuation in the
kernel exponent.  Write \(A_p(s)=\sum_{h=0}^p a_{p,h}s^h\).
For \(1\le r\le p\), define
\begin{equation}\label{bhr:cxp:eq:Delta}
 \Delta_{m,p,r}(z)=m!
 \sum_{h=0}^{\min(p,m)}
 \frac{(-1)^{m-h}a_{p,h}}{(m-h+1)!}\,
 \left.\partial_\lambda^{\,m-h+1}
                 D_{p-r}(\lambda,z)\right|_{\lambda=r}.
\end{equation}
Then
\begin{equation}\label{bhr:cxp:eq:lambdaCT}
 \operatorname{FP}_x\int_0^1
   \gamma_m^{(p)}(x)R_z(x)^r\,dx
 =\operatorname{CT}_{\lambda=r}
      \mathcal M_{m,p}(\lambda,z)-\Delta_{m,p,r}(z).
\end{equation}
For \(\Re\lambda>0\) away from the positive integers
\(1,\ldots,p\), the meromorphic value of
\(\mathcal M_{m,p}\) equals the pointwise \(x\)-finite part.
\end{theorem}

\begin{proof}
Subtract a sufficiently long full local Taylor jet, so that the
remaining endpoint integral converges normally near \((0,r)\).
Among the explicitly integrated Taylor terms, the only one whose
exponent crosses \(-1\) there is
\[
 A_p(s)D_{p-r}(\lambda,z)
                 x^{s+\lambda-r-1}.
\]
Use a smooth cutoff on a small interval next to zero.
The continued integrals of all other Taylor terms, and the normally
convergent remainder, are holomorphic near \((0,r)\).
The distinguished term has meromorphic integral with polar part
\(A_p(s)D_{p-r}(\lambda,z)/(s+\lambda-r)\).
Replacing the cutoff by one only changes a holomorphic function;
this proves \eqref{bhr:cxp:eq:amplitude} with the indicated full
amplitude, not just its value on the polar hyperplane.

At fixed \(\lambda=r\), the analytically continued integral of
the resonant local term is \(A_p(s)D_{p-r}(r,z)/s\).
Every coefficient of \(x^{s-1}\) is a power of \(\log x\)
divided by \(x\), whose finite part on \((0,1)\) is zero.
Deleting the entire amplitude before expansion proves
\eqref{bhr:cxp:eq:fixed-r}.  Terms on the cutoff complement
cancel the artificial splitting point.

For the second formula, put \(\delta=\lambda-r\) and take
the \(s^m\) coefficient of the polar term in
\eqref{bhr:cxp:eq:amplitude}.  It is
\[
 m!\sum_{h=0}^{\min(p,m)}
  (-1)^{m-h}a_{p,h}
       \frac{D_{p-r}(r+\delta,z)}{\delta^{m-h+1}}.
\]
Its constant Laurent coefficient is exactly
\eqref{bhr:cxp:eq:Delta}.  The holomorphic remainder contributes
the same value along either order of expansion.  This proves
\eqref{bhr:cxp:eq:lambdaCT}.  Away from a crossing, endpoint
subtractions have nonzero exponents, and coefficient extraction
commutes with finite-part integration.
\end{proof}

The theorem supplies the full principal part as well: expand the
polynomial \(D_{p-r}(r+\delta,z)\) in the last display and retain
negative powers of \(\delta\).  Some potential poles cancel.
In particular, \(\Delta_{m,p,p}=0\), since \(D_0=1\);
a collision does not automatically imply a nonzero conversion term.

\paragraph{A nonzero conversion for the simple resolvent.}
The simplest example is \(m=0,p=2,r=1\).  Since
\(A_2(s)=s^2-3s+2\) and \(D_1(\lambda,z)=\lambda z\),
\[
 \Delta_{0,2,1}(z)=2z.
\]
At \(z=i\pi\), set \(\ell=\log(2\pi)-i\pi/2\).
The finite Fourier series
\(R_{i\pi}(x)=(e^{2\pi ix}-1)/(2\pi i)\) gives
\[
 \operatorname{FP}_x\int_0^1\gamma_0''(x)R_{i\pi}(x)\,dx
       =2\pi i\left(\frac32-\gamma-\ell\right).
\]
It follows that
\begin{equation}\label{bhr:cxp:eq:noncommute}
 \operatorname{CT}_{\lambda=1}
       \mathcal M_{0,2}(\lambda,i\pi)
       =2\pi i\left(\frac52-\gamma-\ell\right).
\end{equation}
The two constants differ by \(2\pi i\).  This is a concrete
obstruction to identifying exponent continuation with the
specified \(x\)-finite part.  In contrast the case
\(m=0,p=r=1\) has no exponent-axis correction.

\paragraph{Exact recovery of integral powers.}
For a positive integer \(r\) and \(q\ne0\), partial fractions give
\[
 h_n(r,q)=q^{n-r}
   \sum_{j=1}^r\binom rj(q-1)^j
                    \binom{n+j-1}{j-1}\qquad(n\ge1).
\]
The binomial on the right is a polynomial of degree \(j-1\) in
\(n\).  Consequently \(\mathcal L_r(u,q)\) is a finite linear
combination of the ordinary \(\operatorname{Li}_{u-k}(q)\),
\(0\le k<r\), with rational coefficients in \(q\).
At \(q=0\) the removable limiting expression is
\[
 \mathcal L_r(u,0)
       =\sum_{n=1}^r\frac{(-1)^n\binom rn}{n^u}.
\]
Thus the completed generator recovers precisely the ordinary
polylogarithm class of the previous repeated-resolvent theorem.

\subsection{Gamma closure at a special point}

For \(z=\epsilon i\pi\) one has \(q=0\), and
\[
 \mathcal L_\lambda(u,0)
       =\sum_{n\ge1}\frac{(-\lambda)_n}{n!\,n^u}.
\]
At every positive integer \(u\), this new coordinate reduces to
generalized harmonic numbers at a complex argument.  Define
\[
 H_\lambda^{(1)}=\psi(1+\lambda)+\gamma,\qquad
 H_\lambda^{(r)}=\zeta(r)-\zeta(r,1+\lambda)\quad(r\ge2),
\]
initially for \(\Re\lambda>-1\) and then meromorphically.

\begin{theorem}[Every integral spectral order at \(q=0\)]
\label{bhr:cxp:thm:gamma}
For \(N\ge1\),
\begin{equation}\label{bhr:cxp:eq:Bell}
 \mathcal L_\lambda(N,0)
 =-\frac1{N!}Y_N\bigl(
     H_\lambda^{(1)},1!H_\lambda^{(2)},\ldots,
                          (N-1)!H_\lambda^{(N)}\bigr),
\end{equation}
where \(Y_N\) is the complete exponential Bell polynomial.
Equivalently, near \(v=0\),
\begin{equation}\label{bhr:cxp:eq:beta-generator}
 \sum_{N\ge1}\mathcal L_\lambda(N,0)v^{N-1}
 =\frac1v\left\{
 1-\frac{\Gamma(1-v)\Gamma(1+\lambda)}
                 {\Gamma(1+\lambda-v)}\right\}.
\end{equation}
In particular,
\[
 \mathcal L_\lambda(1,0)=-H_\lambda^{(1)},\qquad
 \mathcal L_\lambda(2,0)
 =-\frac{(H_\lambda^{(1)})^2+H_\lambda^{(2)}}2.
\]
\end{theorem}

\begin{proof}
For \(\Re\lambda>-1\) and \(\Re t>0\), the beta integral gives
\[
 \sum_{n\ge1}\frac{(-\lambda)_n}{n!(n+t)}
  =\int_0^1 y^{t-1}\bigl((1-y)^\lambda-1\bigr)\,dy
  =\frac{\Gamma(t)\Gamma(1+\lambda)}
           {\Gamma(1+\lambda+t)}-\frac1t.
\]
Expand at \(t=-v\).  The left side generates the positive
integral spectral orders.  On the right,
\[
 \log\frac{\Gamma(1-v)\Gamma(1+\lambda)}
              {\Gamma(1+\lambda-v)}
   =\sum_{r\ge1}\frac{H_\lambda^{(r)}}r\,v^r.
\]
The defining Bell-polynomial generating function proves both
identities.  Meromorphic continuation removes the initial
restrictions.  This is a beta-integral specialization, not a claim
of priority for the classical binomial-harmonic identity.
\end{proof}

\subsection{Exponent jets and colored multiple polylogarithms}

There is also an exact finite description of every exponent jet at
\(\lambda=0\).  For \(q\ne0,\ |q|<1\), let
\(\mathcal A=\{1,q\}\), with signs \(\nu(1)=-1,\nu(q)=1\).
The multiple polylogarithm uses decreasing indices:
\[
 \operatorname{Li}_{s_1,\ldots,s_d}(z_1,\ldots,z_d)
 =\sum_{n_1>\cdots>n_d\ge1}
       \frac{z_1^{n_1}\cdots z_d^{n_d}}
            {n_1^{s_1}\cdots n_d^{s_d}}.
\]

\begin{proposition}[All exponent jets at zero]
\label{bhr:cxp:prop:jets}
For \(d\ge1\) and \(\Re u>0\),
\begin{equation}\label{bhr:cxp:eq:MPL}
 \left.\partial_\lambda^d\mathcal L_\lambda(u,q)
                                      \right|_{\lambda=0}
 =d!\sum_{a_1,\ldots,a_d\in\mathcal A}
 \left(\prod_{j=1}^d\nu(a_j)\right)
 \operatorname{Li}_{u+1,\{1\}^{d-1}}
       \left(a_1,\frac{a_2}{a_1},\ldots,
                                \frac{a_d}{a_{d-1}}\right).
\end{equation}
All displayed nested sums converge absolutely.  Their cumulative
color products are \(a_j\in\{1,q\}\), so the appearance of
\(q^{-1}\) as an individual color causes no convergence problem.
For nonintegral \(u\), these are nested Dirichlet sums with a complex first index; at positive integral \(u\), they are classical colored multiple polylogarithms.
At \(q=0\), the continuous limiting identity is
\[
 \left.\partial_\lambda^d\mathcal L_\lambda(u,0)
                                      \right|_{\lambda=0}
       =(-1)^d d!\,\zeta(u+1,\{1\}^{d-1}).
\]
\end{proposition}

\begin{proof}
Put \(F(w)=\operatorname{Li}_1(qw)-\operatorname{Li}_1(w)\).
The \(d\)-th exponent derivative of the disk kernel at zero is
\(F(w)^d\).  Repeated integration of
\[
 F'(w)=\frac{q}{1-qw}-\frac1{1-w}
\]
proves the shuffle identity
\[
 \frac{F(w)^d}{d!}
 =\sum_{a_1,\ldots,a_d\in\mathcal A}
   \left(\prod_j\nu(a_j)\right)
 \operatorname{Li}_{\{1\}^d}
     \left(a_1w,\frac{a_2}{a_1},\ldots,
                                   \frac{a_d}{a_{d-1}}\right).
\]
Multiplying the coefficient of \(w^n\) by \(n^{-u}\) raises
the first index from one to \(u+1\).  Absolute convergence for
\(\Re u>0\) justifies evaluation at \(w=1\) and all coefficient
interchanges.  If \(q=0\), only the letter one survives.
\end{proof}

The first jet is the simple identity
\[
 \left.\partial_\lambda\mathcal L_\lambda(u,q)\right|_0
       =\operatorname{Li}_{u+1}(q)-\zeta(u+1).
\]
Higher jets explain how generalized harmonic sums and increasing
multiple-polylogarithm depth enter complex powers.  The finite
description is proved for each fixed jet; it does not imply finite
ordinary-polylogarithm closure after summing all jets.

\paragraph{Verification.}
The accompanying program independently compares the spectral
generator with ordinary \(x\)-quadrature at nonintegral complex
parameters, checks the Gamma specialization against Mellin
quadrature, and compares a nearby-exponent Laurent extrapolation
with \eqref{bhr:cxp:eq:noncommute}.  Its recorded numerical residuals
are diagnostics.  The continuation, pole classification, and
all-order conversion rest on the local subtractions above.
